\section{An actual ray-class family and its uniform supply of norm primes}
\label{sec:ray-prime-estimates}

\status{QR1--QR9 construct the literal arithmetic \texttt{LFam} left open
in QW11. QR10--QR24 prove its analytic short-interval supplier, with a
polynomial real-zero gap and constants uniform over every such family.
QR25 combines these results with the already proved QW reductions.
These are proofs of specific arithmetic and analytic suppliers. The
simultaneous geometric choices in the four global branches remain
separate. No Lean compilation or kernel replay is asserted.}

We write $\mathcal B$ for the conductor bound, to distinguish it from
the scale $B=e^{x^\epsilon}$ in QW1--QW9. Groups of ray classes are
written multiplicatively in the arithmetic construction and additively
in \texttt{LFam}; their identity is therefore denoted $1$ or $0$,
respectively. All implied constants below are absolute unless their
dependence on the fixed zero-free width $c_0$ is specified. Existence
of such constants is sufficient for the literal source statements.
The estimates do not use RH or GRH.

\paragraph{QR1: remove the unit ambiguity in the actual ray coset.}
Let $d\ge100$ be prime with $d\equiv3\pmod4$, let $M\ge3$, and set
$K=\mathbb Q(\sqrt{-d})$, $\omega=(1+\sqrt{-d})/2$.
Then $\mathcal O_K=\mathbb Z[\omega]$ and $\operatorname{disc}K=-d$.
Use the actual ray group modulo $M\mathcal O_K$. The construction in
IG1 chooses, for each ray class $g$, an integral ideal $\mathfrak a$
prime to $M$ in the inverse class and $\alpha_0\in\mathfrak a$ with
$\alpha_0\equiv1\pmod {M\mathcal O_K}$. It identifies the ideals in
class $g$ with elements of the coset
$\alpha_0+M\mathfrak a$, modulo units congruent to one modulo $M$,
by $\alpha\mapsto(\alpha)\mathfrak a^{-1}$. The congruence conditions
on $\alpha$ ensure that this quotient is an integral ideal prime to
$M$. The inverse follows from the definition of the ray relation.

In this case the unit quotient is trivial. If a unit is
$u=1+Mt\ne1$, then $|u|=1$ under the complex embedding, whereas the
nonzero algebraic integer $t$ has $|t|^2=N_{K/\mathbb Q}(t)\ge1$.
Thus $M\le M|t|=|u-1|\le2$, a contradiction. Zero is not in the
coset, since every element is congruent to one modulo $M$.
Consequently the preceding map is a bijection, and
\[
 N\big((\alpha)\mathfrak a^{-1}\big)
       =\frac{|\alpha|^2}{N\mathfrak a},\qquad
 \operatorname{covol}(M\mathfrak a)
       =\frac{M^2N\mathfrak a\sqrt d}{2}.                 \tag{QR.1}
\]
These are the original ray class, norm and Euclidean measure; there is
no unit multiplicity to absorb in a later leading constant.

\paragraph{QR2: obtain a reduced binary form with the required bounds.}
Put $N=N\mathfrak a$ and $\Delta=\operatorname{covol}(M\mathfrak a)$.
A shortest nonzero lattice vector $v$ is primitive: otherwise dividing
by an integer of absolute value at least two gives a shorter vector.
It therefore extends to a $\mathbb Z$-basis $v,w$. Subtract an integer
multiple of $v$ from $w$ so its projection onto $v$ has absolute value
at most $|v|/2$. Since $|w|\ge|v|$, its perpendicular height is at
least $\sqrt3|v|/2$. Hence
\[
 |v|^2\le2\Delta/\sqrt3\le2\Delta.
\]
Also $v=Mz$ with $z\in\mathfrak a\setminus\{0\}$, so
$(z)\subseteq\mathfrak a$ and $|v|^2\ge M^2N$.
Writing the squared norm in this basis gives exactly
\[
 Q(m,n)=\frac{|\alpha_0+mv+nw|^2}{N}
       =a(m+\theta_1+\ell n)^2+b(n+\theta_2)^2,           \tag{QR.2}
\]
where
\[
 a=|v|^2/N,\quad b=\Delta^2/(|v|^2N),\quad
 \ell=\frac{\operatorname{Re}(\bar v w)}{|v|^2},\quad
 \theta_1=\frac{\operatorname{Re}(\bar v\alpha_0)}{|v|^2},\quad
 \theta_2=\frac{\operatorname{Im}(\bar v\alpha_0)}
                  {\operatorname{Im}(\bar v w)}.
\]
Orient the basis so that the last denominator is positive; its sign
would not affect (QR.2). The bounds above imply
\[
 a\ge1,\qquad b\ge\frac{\Delta}{2N}
                 =\frac{M^2\sqrt d}{4}\ge1,\qquad
 ab=\frac{dM^4}{4}.                                      \tag{QR.3}
\]
By QR1, every $Q(m,n)$ is a positive integer, and the number of its
representations of $j$ is exactly $a_g(j)$, the number of integral
ideals of norm $j$, prime to $M$, in the original class $g$.

\paragraph{QR3: the ellipse count is uniform in its shifts and shear.}
For $x\ge1$ and fixed $n$, the condition $Q(m,n)\le x$ is an interval
in $m$ of length
\[
 f(n+\theta_2),\qquad
 f(t)=\frac2{\sqrt a}\sqrt{(x-bt^2)_+}.
\]
Its integer count differs from its length by at most one. There are
at most $2\sqrt{x/b}+3$ rows meeting the ellipse. For any even,
nonnegative function decreasing on $[0,\infty)$ with compact support,
\[
 \left|\sum_{n\in\mathbb Z}f(n+\theta)-\int_{\mathbb R}f(t)\,dt\right|
 \le f(0).                                                \tag{QR.4}
\]
Indeed write $f$ as the integral of the indicators of its superlevel
sets. Each nonempty superlevel set is a centered interval, and its
shifted integer count differs from its length by at most one.
Integrating that discrepancy over the levels proves (QR.4); endpoints
of levels do not affect the integral. The semicircle integral is
$\pi x/\sqrt{ab}$. Thus (QR.3) gives
\[
 \left|\#\{(m,n):Q(m,n)\le x\}-\frac{\pi x}{\sqrt{ab}}\right|
 \le2\sqrt{x/b}+3+2\sqrt{x/a}\le10\sqrt x.                 \tag{QR.5}
\]
The row intervals also prove finiteness directly. Neither $\ell$ nor
the two shifts changes this estimate.

\paragraph{QR4: a uniform shifted Gaussian estimate.}
For $s>0$, periodize $e^{-s t^2}$ on $\mathbb R/\mathbb Z$.
Its Fourier coefficients are
$\sqrt{\pi/s}\,e^{-\pi^2j^2/s}$. The Gaussian and the resulting
Fourier series converge absolutely and uniformly on a period, so
integration against the characters and Fourier inversion give
\[
 \sum_{m\in\mathbb Z}e^{-s(m+\theta)^2}
 =\sqrt{\pi/s}\sum_{j\in\mathbb Z}
                     e^{-\pi^2j^2/s}e^{2\pi i j\theta}.   \tag{QR.6}
\]
The Gaussian Fourier transform follows by completing the square, or
by its differential equation and its value at zero. The periodization
argument is the same elementary Poisson argument used in IG1.
Since $e^{-y}\le2y^{-2}$ for $y>0$ and
$\sum_{j\ge1}j^{-4}\le2$,
\[
 2\sum_{j\ge1}e^{-\pi^2j^2/s}
 \le\frac{4s^2}{\pi^4}\sum_{j\ge1}j^{-4}\le s^2.
\]
Subtracting the term $j=0$ in (QR.6) therefore proves, for every shift,
\[
 \left|\sum_m e^{-s(m+\theta)^2}-\sqrt{\pi/s}\right|
                  \le\sqrt{\pi/s}\,s^2.                 \tag{QR.7}
\]
All sums involved are absolutely convergent, so the same estimate
can be used inside a second Gaussian sum.

\paragraph{QR5: verify both actual ray counting laws.}
Apply QR4 to the inner sum in $m$ in $e^{-tQ(m,n)}$; its shift
$\theta_1+\ell n$ depends on $n$, but the error bound does not.
Apply it again to the remaining sum in $n$. Bounding that outer
sum by its Gaussian main term times $1+(bt)^2$ gives
\[
 \left|\sum_{m,n}e^{-tQ(m,n)}-\frac{\pi}{\sqrt{ab}\,t}\right|
 \le\frac{\pi}{\sqrt{ab}\,t}
           \big((at)^2+(bt)^2+(at)^2(bt)^2\big).            \tag{QR.8}
\]
For $0<t\le1$, (QR.3) bounds this by $3\pi(ab)^2t$.
The equality between the double sum and
$\sum_{j\ge0}a_g(j)e^{-tj}$ follows by partition into the finite
norm fibers of QR2; nonnegativity also permits Tonelli's theorem.
Set
\[
 \mathcal B=4dM^2,\qquad
 \kappa=\frac{2\pi}{\sqrt d\,M^2}=\frac{\pi}{\sqrt{ab}}.
                                                               \tag{QR.9}
\]
The elementary inequalities $d\ge100$, $M\ge3$ and $\pi<4$
give $3\pi(ab)^2\le\mathcal B^{10}$. Combining (QR.5), (QR.8)
and the integral values of $Q$ proves exactly
\[
 \left|\sum_{1\le j\le\lfloor x\rfloor}a_g(j)-\kappa x\right|
 \le10\sqrt x,\qquad
 \left|\sum_{j\ge0}a_g(j)e^{-tj}-\kappa/t\right|
 \le\mathcal B^{10}t.                                    \tag{QR.10}
\]
The first estimate holds for every $x\ge1$; the second for every
$0<t\le1$, uniformly over all ray classes of this same field and
modulus.

\paragraph{QR6: ideal factorization gives the literal convolution.}
Define $\lambda(\mathfrak p^r)=\log N\mathfrak p$ for a nonzero
prime ideal $\mathfrak p$ and $r\ge1$, and define it to be zero on
other nonzero ideals. Unique ideal factorization gives
\[
 \sum_{\mathfrak j\mid\mathfrak i}\lambda(\mathfrak j)
       =\sum_{\mathfrak p}v_{\mathfrak p}(\mathfrak i)
                                \log N\mathfrak p
       =\log N\mathfrak i.                               \tag{QR.11}
\]
Let $\lambda_g(n)$ be the sum of these weights over the ideals counted
by $a_g(n)$. The norm fibers and the ideal divisor sets are finite.
For an ideal prime to $M$, both factors in a factorization
$\mathfrak i=\mathfrak j\mathfrak k$ are prime to $M$; conversely
the product of two such ideals is prime to $M$. Their norms multiply
and their ray classes add. Sum (QR.11) over the fiber of $(g,n)$ and
use this factorization bijection to obtain
\[
 a_g(n)\log n
  =\sum_{h\in G}\ \sum_{uv=n}\lambda_h(u)a_{g-h}(v).       \tag{QR.12}
\]
For $n=0$ the two sides are zero under the source convention, and
for $n=1$ every ideal weight is zero. No rearrangement of a conditionally
convergent series occurs here.

\paragraph{QR7: the degree-two weight bound retains norm-one support.}
All the weights are nonnegative. A nonzero weight at an integer norm
is supported on $n=p^r$. If $N\mathfrak p=p^f$, the contribution
of $\mathfrak p^k$ at that norm occurs only when $fk=r$ and equals
$f\log p$. For fixed $p$, the identity
$\sum_{\mathfrak p\mid p}e_{\mathfrak p}f_{\mathfrak p}=2$
implies $\sum f_{\mathfrak p}\le2$. Removing ideals not prime to
$M$ can only decrease the sum. Hence
\[
             \sum_g\lambda_g(n)\le2\Lambda(n).            \tag{QR.13}
\]
If $p$ is prime and $\lambda_0(p)\ne0$, there is an integral ideal
$\mathfrak i$ of norm $p$ in the identity ray class. The defining
ray relation has $(a)\mathfrak i=(b)$ with nonzero $a$ and
$a,b\equiv1\pmod M$. Since $(b)\subseteq(a)$, write $b=a\pi$
with $\pi\in\mathcal O_K$. Cancellation gives $\mathfrak i=(\pi)$,
and
$\pi-1=(a\pi-1)-\pi(a-1)\in M\mathcal O_K$.
Write $\pi=1+M(r+s\omega)$, $r,s\in\mathbb Z$, and put
$A=2(1+Mr)+Ms$, $D=Ms$. The exact quadratic norm identity gives
\[
       4p=A^2+dD^2,\qquad M\mid D,\qquad
                          2M\mid A-D-2.                  \tag{QR.14}
\]
These are precisely the original \texttt{NormOne} conditions, with
the same $d,M,p$ as the family. In particular no unsupported
replacement of a ray condition by ordinary principalness has been made.

\paragraph{QR8: real-character convolution is nonnegative and counts squares.}
Let $\chi^2=1$ be a ray character, so $\chi$ takes values $\pm1$.
For an ideal prime to $M$ put
\[
 r_\chi(\mathfrak i)=\sum_{\mathfrak j\mid\mathfrak i}
                  \chi([\mathfrak j])
 =\prod_{\mathfrak p^e\Vert\mathfrak i}
       (1+\chi([\mathfrak p])+\cdots+\chi([\mathfrak p])^e).
                                                               \tag{QR.15}
\]
If the sign is $1$, the factor is $e+1$; if it is $-1$, the factor
is $0$ for odd $e$ and $1$ for even $e$. All factors are nonnegative,
and every factor of $r_\chi(\mathfrak j^2)$ is at least one.
Grouping ideal factorizations by their product shows that the ordinary
Dirichlet convolution coefficient
\[
 c_\chi(n)=\sum_{uv=n}\Big(\sum_g a_g(u)\Big)
                              \Big(\sum_g\chi(g)a_g(v)\Big)
       =\sum_{N\mathfrak i=n}r_\chi(\mathfrak i)           \tag{QR.16}
\]
is nonnegative; on the right, all ideals are prime to $M$.
Squaring such ideals is injective by ideal factorization. Keeping just
the terms $\mathfrak i=\mathfrak j^2$ in $c_\chi(n^2)$ proves
$c_\chi(n^2)\ge\sum_g a_g(n)$. The zero norm case is also zero.
This proves both clauses of the source \texttt{sq} field, not merely
positivity of a few initial coefficients.

\paragraph{QR9: assemble every field of the actual arithmetic family.}
The ideal of norm one is $\mathcal O_K$, in the identity class;
there is no nonzero ideal of norm zero. Thus
$a_0(1)=1$, $a_g(1)=0$ for $g\ne0$, and $a_g(0)=0$.
The ray exact sequence bounds
$|G|\le h_K\,|\mathcal O_K/M\mathcal O_K|=h_KM^2$.
Here the short-vector argument of TL14 has the following explicit
two-dimensional specialization. An integral ideal $I$ has covolume
$\Delta_I=N(I)\sqrt d/2$. A disk of radius $\sqrt{\Delta_I}$ has
area $\pi\Delta_I>\Delta_I$. Averaging its lattice-point count over
a fundamental domain therefore gives a translate containing two
distinct points of $I$. Their nonzero difference $a$ satisfies
$|a|^2\le4\Delta_I=2\sqrt d\,N(I)$. Thus the integral ideal
$(a)I^{-1}$ has norm at most $2\sqrt d$; choosing $I$ in the inverse
class supplies such a representative in every class.
The number of degree-two ideals of
norm $n$ is at most the divisor function $\tau(n)$, by considering
split, inert and ramified prime powers. Therefore the number with
norm at most $2\sqrt d$ is at most
$\lfloor2\sqrt d\rfloor^2\le4d$. Hence $|G|\le\mathcal B$.
Formula (QR.9) gives $\mathcal B\ge4$ and
$1/\mathcal B\le\kappa\le4$ directly.

Take the actual finite additive ray group, $\mathcal B,\kappa$, and
the actual coefficients $a_g,\lambda_g$. Their nonnegativity and
norm-zero and norm-one conditions have just been proved. QR5 gives
both counting fields; QR6 gives convolution; QR7 gives the total
von Mangoldt bound; QR8 gives both real-character clauses. These
are all fourteen laws in the \texttt{LFam} record. QR7 simultaneously
supplies its norm-one support. This constructs
\texttt{exists\_lFam\_ray} and its public statement with the exact
conductor $4dM^2$; no arithmetic \texttt{LFam} supplier remains
as a hypothesis of QW11.

\paragraph{QR10: continue the whole family by its actual counting remainder.}
For the rest of this section let $\mathcal F$ be any \texttt{LFam},
with group $G$ of order $h$. Put
\[
 a_\chi(n)=\sum_g\chi(g)a_g(n),\quad
 \lambda_\chi(n)=\sum_g\chi(g)\lambda_g(n),\quad
 \rho_\chi=\begin{cases}\kappa h,&\chi=1,\\0,&\chi\ne1.\end{cases}
\]
Finite character orthogonality and the count field imply
\[
 R_\chi(x)=\sum_{n\le x}a_\chi(n)-\rho_\chi x,
       \qquad |R_\chi(x)|\le10\mathcal B\sqrt x
                      \quad(x\ge1).                       \tag{QR.17}
\]
Consequently
$m_\chi(s)=\int_1^\infty R_\chi(x)x^{-s-1}\,dx$
is holomorphic for $\operatorname{Re}s>1/2$, with
$|m_\chi(s)|\le10\mathcal B/(\operatorname{Re}s-1/2)$.
On every compact sub-half-plane its differentiated integrands are
dominated by $C x^{-1-\eta}(1+\log x)^j$; this proves both convergence
and holomorphy by differentiation under the integral.
Abel summation, justified first for $\operatorname{Re}s>1$, gives
\[
 L_\chi(s):=\sum_{n\ge1}a_\chi(n)n^{-s}
       =\frac{\rho_\chi s}{s-1}+s m_\chi(s).              \tag{QR.18}
\]
Thus the literal regularized function is
\[
 F_\chi(s)=\begin{cases}
 \rho_\chi s+(s-1)s m_\chi(s),&\chi=1,\\
 s m_\chi(s),&\chi\ne1.
 \end{cases}                                               \tag{QR.19}
\]
It is holomorphic for $\operatorname{Re}s>1/2$, equals
$(s-1)L_1(s)$ or $L_\chi(s)$ in the absolutely convergent region,
and satisfies
\[
 |F_\chi(s)|\le300\mathcal B(|s|+1)^2
                    \quad(\operatorname{Re}s\ge11/20).    \tag{QR.20}
\]
Indeed $|\rho_\chi|\le4\mathcal B$ and the preceding Mellin bound
is at most $200\mathcal B$ there. Real characters give
$F_\chi(\bar s)=\overline{F_\chi(s)}$ by conjugating the integral.

\paragraph{QR11: prove nonvanishing and the logarithmic derivative to the right.}
The nonnegative weights and their total bound give
$|\lambda_\chi(n)|\le2\Lambda(n)$ and $\lambda_\chi(1)=0$.
Take the finite character transform of (QR.12). Absolute convergence
for $\sigma>1$ permits multiplication of the resulting Dirichlet
series and gives
\[
 L'_\chi(s)=-G_\chi(s)L_\chi(s),\qquad
 G_\chi(s)=\sum_{n\ge2}\lambda_\chi(n)n^{-s}.             \tag{QR.21}
\]
Absolute convergence of $L_\chi$ in this region follows from its
partial-sum bound $\sum_{n\le x}|a_\chi(n)|\le C\mathcal Bx$;
termwise derivatives follow on smaller half-planes. Define
\[
 H_\chi(s)=\sum_{n\ge2}\frac{\lambda_\chi(n)}{\log n}\,n^{-s}.
\]
It has derivative $-G_\chi(s)$. Therefore
$L_\chi(s)e^{-H_\chi(s)}$ has derivative zero on the connected
half-plane. Its limit along the positive real axis is one, since
$a_\chi(1)=1$ and all remaining absolutely convergent terms tend to
zero. It is identically one. In particular
\[
 L_\chi=e^{H_\chi}\ne0,\qquad
 G_\chi(s)=\frac{\mathbf1_{\chi=1}}{s-1}
                        -\frac{F'_\chi(s)}{F_\chi(s)}.    \tag{QR.22}
\]
Unique integer factorization yields the absolutely convergent identity
$\sum_{n\ge2}\Lambda(n)n^{-\sigma}/\log n=\log\zeta(\sigma)$.
The integral test gives $\zeta(\sigma)\le\sigma/(\sigma-1)$.
Hence
\[
 |L_\chi(\sigma+it)|\ge e^{-2\log\zeta(\sigma)}
                       \ge((\sigma-1)/\sigma)^2.          \tag{QR.23}
\]
This proves all the required right-half-plane assertions without
presupposing an Euler product for the abstract family.

\paragraph{QR12: turn the theta field into an integrable nonprincipal heat sum.}
Put $\Theta_\chi(t)=\sum_n a_\chi(n)e^{-tn}$ for $t>0$.
From the count bound and nonnegativity,
$a_g(n)\le4n+10\sqrt n\le14n$ for $n\ge1$; consequently
$|a_\chi(n)|\le14\mathcal B n$. These bounds prove absolute
convergence of every heat sum and justify the finite character sum.
For $0<t\le1$, the theta field and orthogonality give
\[
       |\Theta_\chi(t)-\rho_\chi/t|\le\mathcal B^{11}t.
\]
For $t\ge1$, summing $n e^{-tn}$ bounds $|\Theta_\chi(t)|$
by $C\mathcal B e^{-t}$. After increasing an absolute $C\ge1$,
these statements imply
\[
 |\Theta_\chi(t)-\rho_\chi/t|
      \le C\mathcal B^{11}\min(t,t^{-1})\quad(t>0).       \tag{QR.24}
\]
If $\chi\ne1$, the small-$t$ bound and the stronger exponential
large-$t$ bound make $\Theta_\chi$ integrable. Its integral is
exactly $F_\chi(1)$. To see this without exchanging a possibly
nonabsolute series at $s=1$, use Abel summation in the form
\[
 \Theta_\chi(t)=t\int_1^\infty A_\chi(x)e^{-tx}\,dx,
      \qquad A_\chi(x)=\sum_{n\le x}a_\chi(n).
\]
For nonprincipal $\chi$, (QR.17) bounds $|A_\chi(x)|$ by
$10\mathcal B\sqrt x$ and
$\int_1^\infty|A_\chi(x)|x^{-2}\,dx<\infty$.
Fubini and $\int_0^\infty t e^{-tx}\,dt=x^{-2}$ now give
\[
                  \int_0^\infty\Theta_\chi(t)\,dt
                          =m_\chi(1)=F_\chi(1).           \tag{QR.25}
\]

\paragraph{QR13: a positive Bessel integral gives a polynomial lower bound at one.}
Let $\chi\ne1$ be real. Write $P=\Theta_1$, $Q=\Theta_\chi$,
$\rho=\kappa h$, $E=C\mathcal B^{11}$. The \texttt{sq} field
says the Dirichlet convolution coefficients $c_\chi(n)$ of
$a_1$ and $a_\chi$ are nonnegative. The norm-one coefficients give
$c_\chi(1)=1$. Define the positive kernel
\[
 W(y)=\int_0^\infty e^{-y/w-w}\,\frac{dw}{w}\quad(y>0),
       \qquad w_0=W(1)>0.
\]
The inequality $e^{-y/w}\le6w^3/y^3$ proves $W(y)\le12y^{-3}$.
Together with $|a_1(n)|,|a_\chi(n)|\le14\mathcal B n$, this
shows absolute integrability of the expanded double series below:
its absolute integrals are bounded by a constant times
$\mathcal B^2v^{-3}\sum_{m,n\ge1}m^{-2}n^{-2}$.
Substitute $w\mapsto w/m$ in each term to obtain
\[
 I(v):=\int_0^\infty P(v/w)Q(w)\,\frac{dw}{w}
       =\sum_n c_\chi(n)W(nv)\ge W(v)\ge w_0
                                      \quad(0<v\le1).     \tag{QR.26}
\]
Since $\min(t,t^{-1})\le\sqrt t$, (QR.24) gives
\[
 \left|I(v)-\frac\rho v\int_0^\infty Q(w)\,dw\right|
 \le E^2\sqrt v\,K_0,\qquad
 K_0=\int_0^\infty\frac{\min(w,w^{-1})}{w\sqrt w}\,dw
                    =2+\frac23=\frac83.                  \tag{QR.27}
\]
The integrable majorant in (QR.27) justifies subtracting the main
term. Choose once and for all
\[
 \delta=\min\left(1,\frac{w_0}{2(C^2K_0+1)}\right),\quad
 s=\delta\mathcal B^{-22},\quad v=s^2,\quad
 c_* =\min(1,\delta^2w_0/8)>0.
\]
The error in (QR.27) is at most $w_0/2$, so (QR.25)--(QR.27) and
$\rho\le4\mathcal B$ imply
$|F_\chi(1)|\ge w_0s^2/(8\mathcal B)\ge c_*\mathcal B^{-45}$.
For the principal character, (QR.19) gives
$F_1(1)=\rho\ge1/\mathcal B$, which satisfies the same lower bound.
Thus for every real character, including the principal one,
\[
                      |F_\chi(1)|\ge c_*\mathcal B^{-45}. \tag{QR.28}
\]
Every constant was chosen independently of the family. In particular
this step does not invoke an ineffective exceptional-zero theorem.

\paragraph{QR14: the real-zero gap follows from a local derivative bound.}
If $|z-1|\le1/4$, the disk $|w-z|\le1/10$ lies in
$\operatorname{Re}w\ge13/20$ and $|w|\le27/20$.
The growth bound (QR.20) is therefore at most $1800\mathcal B$
on that disk. Cauchy's derivative estimate gives
\[
                   |F'_\chi(z)|\le18000\mathcal B.        \tag{QR.29}
\]
For a real character and a real zero $\beta\ge3/4$, QR11 rules
out $\beta>1$ and QR13 rules out $\beta=1$. Integrating the derivative
along $[\beta,1]$ and using (QR.28) yields
\[
 c_*\mathcal B^{-45}\le|F_\chi(1)-F_\chi(\beta)|
           \le18000\mathcal B(1-\beta),\qquad
 1-\beta\ge c_9\mathcal B^{-46},\quad
 c_9=c_*/18000\in(0,1].                                  \tag{QR.30}
\]
This is the literal uniform real-zero separation needed below. It
is a polynomial separation for these families, with exponent $46$;
it is not a fixed-width exclusion of real zeros for unbounded
conductors.

\paragraph{QR15: a finite Blaschke product controls the local zero sum.}
Here is the complex-analysis estimate used in the zero-free argument.
Fix $0<r<R<1$. Suppose $f$ is holomorphic on a neighborhood of the
closed unit disk, $f(0)\ne0$, and $|f|\le M$ there. Its zeros
$a_j$ in $|z|\le R$, with multiplicities $m_j$, form a finite set:
an accumulation point would force the analytic function to vanish
identically. Divide by the finite Blaschke product
\[
 B_f(z)=\prod_j\left(\frac{z-a_j}{1-\bar a_jz}\right)^{m_j},
                    \qquad g=f/B_f.
\]
Removable singularities at the $a_j$ make $g$ holomorphic on the
unit disk and its boundary. On the unit circle $|B_f|=1$, so the
maximum principle gives $|g|\le M$. At zero,
$|g(0)|=|f(0)|/\prod_j|a_j|^{m_j}$. Thus, with
$L=\log(M/|f(0)|)\ge0$ and $N=\sum_jm_j$,
\[
                     N\log(1/R)\le L.                    \tag{QR.31}
\]
The function $g$ has no zeros on the disk of radius $R$, so an
analytic logarithm exists there. The real part of
$H=\log M-\log g$ is nonnegative, and
$\operatorname{Re}H(0)\le L$. The Poisson formula for this real
part, differentiated in the smaller disk, gives
\[
       |H'(z)|\le\frac{2R}{(R-r)^2}\operatorname{Re}H(0)
                                        \quad(|z|\le r).
\]
One may obtain the displayed bound by integrating the derivative of
$(Re^{i\theta}+z)/(Re^{i\theta}-z)$ against the nonnegative boundary
values; its absolute value is at most $2R/(R-r)^2$, and their mean
is $\operatorname{Re}H(0)$. A radius tending down to $R$ handles
any boundary issue. Finally
\[
 \frac{f'}f-\sum_j\frac{m_j}{z-a_j}
       =\frac{g'}g+\sum_j\frac{m_j\bar a_j}{1-\bar a_jz}.
\]
The last sum is at most $N/(1-r)$ in absolute value. With (QR.31)
this proves, for a constant depending only on $r,R$,
\[
 N\le C_{r,R}(1+L),\qquad
 \left|\frac{f'(z)}{f(z)}-\sum_j\frac{m_j}{z-a_j}\right|
       \le C_{r,R}(1+L)\quad(|z|\le r,\ f(z)\ne0).        \tag{QR.32}
\]
This is a direct proof of the local logarithmic-derivative lemma;
no unproved global zero-count estimate is substituted for it.

\paragraph{QR16: specialize that estimate with uniform family constants.}
Put $\ell(t)=\log(\mathcal B(|t|+2))>2$ and center a disk at
$z_t=5/4+it$. Formula (QR.23) gives
$|F_\chi(z_t)|\ge1/100$, including the principal factor $s-1$.
On the disk of radius $7/10$, (QR.20) gives the upper bound
$300\mathcal B(|t|+3)^2$. Rescale this disk to the unit disk
and apply QR15 with $r=5/7$, $R=6/7$. The logarithm of the
upper-to-lower ratio, plus one, is at most a fixed multiple of
$\ell(t)$. There is therefore one absolute $C>0$ such that
the zeros in $|\rho-z_t|\le3/5$, counted with multiplicity,
satisfy
\[
 \sum_\rho m_\rho\le C\ell(t),\qquad
 \left|\frac{F'_\chi(s)}{F_\chi(s)}
                   -\sum_\rho\frac{m_\rho}{s-\rho}\right|
          \le C\ell(t)\quad(|s-z_t|\le1/2,\ F_\chi(s)\ne0).
                                                               \tag{QR.33}
\]
All zeros have real part at most one by QR11. Consequently, at
$s=1+a+it$, $0<a\le1/2$, every term
$\operatorname{Re}(s-\rho)^{-1}$ is nonnegative. Combining
(QR.22) and (QR.33), and keeping any selected zeros in that disk,
gives
\[
 \operatorname{Re}G_\chi(1+a+it)
 \le\mathbf1_{\chi=1}\frac{a}{a^2+t^2}+C\ell(t)
         -\sum_{\rho\text{ selected}}m_\rho
                                      \operatorname{Re}\frac1{s-\rho}.
                                                               \tag{QR.34}
\]
The same $C$ can be increased once to cover all subsequent uses.
The elementary comparisons $\ell(0)\le\ell(t)$ and
$\ell(2t)\le2\ell(t)$ will be used below.

\paragraph{QR17: positivity gives a zero-free region outside real characters.}
For $|z|=1$,
$3+4\operatorname{Re}z+\operatorname{Re}z^2
 =2(1+\operatorname{Re}z)^2\ge0$.
Apply this with $z=\chi(g)n^{-it}$, multiply by
$\lambda_g(n)n^{-\sigma}\ge0$, and sum absolutely for $\sigma>1$.
This proves
\[
 3G_1(\sigma)+4\operatorname{Re}G_\chi(\sigma+it)
              +\operatorname{Re}G_{\chi^2}(\sigma+2it)\ge0.\tag{QR.35}
\]
Also $G_1(\sigma)\ge0$; if $\chi$ is real, then
$G_1(\sigma)+G_\chi(\sigma)\ge0$. For two real characters,
\[
 G_1(\sigma)+G_\chi(\sigma)+G_\psi(\sigma)
                                +G_{\chi\psi}(\sigma)\ge0,\tag{QR.36}
\]
because the corresponding coefficient is
$(1+\chi(g))(1+\psi(g))\ge0$.

We give the numerical separation step to fix the quantifiers.
Suppose $F_\chi(\rho)=0$, $\rho=1-\delta+i\gamma$,
$0\le\delta\le\varepsilon\le1/400$, and either $\chi^2\ne1$
or $\gamma\ne0$. If $\chi$ is real and
$0<|\gamma|<30\varepsilon$, its conjugate is a distinct zero.
Both are in the disk of (QR.33) centered at $5/4$.
At $\sigma=1+60\varepsilon$, their combined contribution is
at least
\[
 \frac{2(60\varepsilon)}{(61\varepsilon)^2+(30\varepsilon)^2}
                         =\frac{120}{4621\varepsilon}.
\]
Use $G_1\ge0$ if $\chi=1$, and $G_1+G_\chi\ge0$ otherwise.
Equation (QR.34), with the possible principal pole kept, yields
\[
 \frac{120}{4621\varepsilon}
          \le\frac1{60\varepsilon}+2C\ell(0),
       \qquad \varepsilon^{-1}\le300C\ell(\gamma).        \tag{QR.37}
\]
In all other cases use (QR.35) at $\sigma=1+6\varepsilon$,
$t=\gamma$, keeping the zero $\rho$ in the middle term.
Its contribution is at least $4/(7\varepsilon)$. If
$\chi^2\ne1$, the shifted pole terms are absent. Otherwise
$|\gamma|\ge30\varepsilon$, and each possible shifted pole
term is at most $1/(156\varepsilon)$ before its coefficient.
The principal term contributes $3/(6\varepsilon)$, and the
total errors are at most $9C\ell(\gamma)$. Hence
\[
 \frac4{7\varepsilon}
   \le\frac1{2\varepsilon}+\frac5{156\varepsilon}
                                  +9C\ell(\gamma),
       \qquad \varepsilon^{-1}\le300C\ell(\gamma).        \tag{QR.38}
\]
Both final inequalities are direct rational comparisons.
Put $K=300C$ and $c_{\rm zf}=1/(4K+400)>0$.
If a zero satisfies
$\operatorname{Re}\rho\ge1-c_{\rm zf}/\ell(\operatorname{Im}\rho)$,
choose $\varepsilon=1/((2K+200)\ell(\gamma))$.
Then $0\le\delta\le\varepsilon/2$, $\varepsilon\le1/400$,
and (QR.37) or (QR.38) is impossible. Every zero in this region
is therefore real and belongs to a real character. This proves
the stated family zero-free region, allowing the precise exception
that the next steps control.

\paragraph{QR18: at most one exceptional zero occurs in the entire family.}
Increase the constant $C$ of (QR.34) if needed, and choose
$c_{\rm un}=1/(100(C+1))$. Put $\ell_0=\ell(0)$ and
$\sigma=1+2c_{\rm un}/\ell_0$.
Any real zero $\beta\ge1-c_{\rm un}/\ell_0$ lies in the
disk used in (QR.34) at $t=0$ and contributes at least
\[
             \frac{m_\beta}{\sigma-\beta}
                    \ge\frac{m_\beta\ell_0}{3c_{\rm un}}.\tag{QR.39}
\]
Suppose that two such zero occurrences exist, counting multiplicity,
for real characters. If they belong to the same nonprincipal
character, use $G_1+G_\chi\ge0$; for the principal character
use $G_1\ge0$. If they belong to distinct nonprincipal real
characters $\chi,\psi$, use (QR.36). The four characters
$1,\chi,\psi,\chi\psi$ are then distinct. If one character
is principal, the corresponding two-term inequality suffices.
In every case exactly one principal pole is present, there are at
most four logarithmic-derivative errors, and (QR.34),(QR.39) force
\[
       0\le\frac{\ell_0}{2c_{\rm un}}+4C\ell_0
                                 -\frac{2\ell_0}{3c_{\rm un}}<0.
\]
The strict inequality follows from $1/(6c_{\rm un})>4C$.
This contradiction proves uniqueness across all characters and
simplicity of the zero if it exists. It also excludes two zeros
of one character; a count of exceptional characters alone would
not have supplied that assertion.

\paragraph{QR19: cancel the exception before estimating the derivative.}
Fix $0<c_0\le1/20$ and a finite set $Z$ satisfying the literal
$\operatorname{ExcZ}(\chi,c_0,Z)$: it contains at most one
real simple zero in $[1-c_0/\ell(0),1)$ and lists every zero
in $\operatorname{Re}s\ge1-c_0/\ell(\operatorname{Im}s)$.
For
\[
 1-\frac{c_0}{2\ell(t)}\le\operatorname{Re}s\le3,
             \quad t=\operatorname{Im}s,\quad s\notin Z,
\]
we claim
\[
 \left|\frac{F'_\chi(s)}{F_\chi(s)}
                         -\sum_{\beta\in Z}\frac1{s-\beta}\right|
                     \le C(c_0)\ell(t)^2.                 \tag{QR.40}
\]
For $\operatorname{Re}s\ge7/4$, use the absolutely convergent
series (QR.22); its absolute bound is uniform, and the real pole
and exception denominators have size at least $3/4$.
Otherwise $s$ lies in the radius-$1/2$ disk of (QR.33).
A zero $\rho$ in the radius-$3/5$ disk has
$\ell(\operatorname{Im}\rho)\le(4/3)\ell(t)$, since its
height differs from $t$ by at most $3/5$. If $\rho\notin Z$,
the defining zero-free condition for $Z$ gives
\[
 \operatorname{Re}(s-\rho)
    \ge\frac{c_0}{\ell(\operatorname{Im}\rho)}
                        -\frac{c_0}{2\ell(t)}
    \ge\frac{c_0}{4\ell(t)}.
\]
There are at most $C\ell(t)$ such zeros with multiplicity;
their reciprocal sum is at most $4C\ell(t)^2/c_0$.
An exceptional zero inside the disk cancels exactly, because it
is simple. If it is outside, its distance from $s$ is at least
$3/5-1/2=1/10$, so its added reciprocal has size at most ten.
Combine these facts with (QR.33) to prove (QR.40). The argument
also proves $F_\chi(s)\ne0$ at the points under consideration.

\paragraph{QR20: truncated Perron with the same integer endpoints.}
Let $b(n)$ be any complex coefficients with $|b(n)|\le2\Lambda(n)$.
For $2\le T\le u$, put $a=1+1/\log u$ and
$D_b(s)=\sum_{n\ge1}b(n)n^{-s}$. There is one absolute $P>0$
such that
\[
 \left|\sum_{1\le n\le\lfloor u\rfloor}b(n)
 -\frac1{2\pi}\int_{-T}^T D_b(a+it)\frac{u^{a+it}}{a+it}\,dt\right|
       \le P\left(\frac{u\log^2u}{T}+\log u\right).       \tag{QR.41}
\]
Here are the kernel and endpoint details. For $y>0$,
\[
 K_T(y)=\frac1{2\pi}\int_{-T}^T\frac{y^{a+it}}{a+it}\,dt.
\]
Closing the infinite vertical integral on the appropriate side gives
$\mathbf1_{y>1}$ when $y\ne1$. Integration by parts in each
tail bounds the error by $C y^a/(T|\log y|)$. A uniform second
bound is $C y^a$: after pairing $t$ and $-t$, the cosine integral
has weight $a/(a^2+t^2)$ of bounded integral, and the sine integral
with weight $t/(a^2+t^2)$ is uniformly bounded by splitting at $a$
and applying the Dirichlet bound to $\sin(t\log y)/t$.
The difference of the latter weight and $1/t$ has bounded integral
on $[a,\infty)$. Thus for $y\ne1$ the error is bounded by
\[
             C y^a\min(1,(T|\log y|)^{-1}).               \tag{QR.42}
\]
At $y=1$, $K_T(1)=\arctan(T/a)/\pi$, so the error from the
full endpoint weight one is at most one.

Absolute convergence on $\operatorname{Re}s=a>1$ permits the
finite-height integral and the coefficient series to be interchanged.
Apply (QR.42) with $y=u/n$. For $u/2<n<2u$,
$|\log(u/n)|\ge|u-n|/(2u)$ and $(u/n)^a$ is bounded, since
$a\le3$. The at most two closest integers cost $O(\log u)$;
the remaining terms cost at most
$C\log u\,(u/T)\sum_{1\le j\le2u}j^{-1}
 =O(u\log^2u/T)$. Outside this interval,
$|\log(u/n)|\ge\log2$ and
\[
 \sum_{n\ge2}\frac{\log n}{n^a}
                   \le C(1+(a-1)^{-2})\le C'\log^2u.
\]
Since $u^a=eu$, this gives the same error for the far terms.
If $u$ is an integer, changing its half endpoint to the required
full weight costs at most $|b(u)|\le2\log u$, already included
in (QR.41). These bounds prove the stated formula for the exact
floor convention, without a hidden distance-to-integers restriction.

\paragraph{QR21: shift the rectangle while keeping the exceptional residue.}
Assume $\operatorname{ExcZ}(\chi,c_0,Z)$ as in QR19, let $u\ge2$
and $T\ge1$, and put
\[
 \tau=\frac{c_0}{8\ell(T)},\qquad
 \mathcal M=C(c_0)\ell(T)^2+2/\tau.
\]
One of the two numbers $1-2\tau$, $1-4\tau$ has distance at
least $\tau$ from the at most one member of $Z$. Choose that
number as $\sigma$. The rectangle from $\sigma-iT$ to $a+iT$,
$a=1+1/\log u\le3$, is contained in the strip of QR19, because
$\ell(t)\le\ell(T)$ for $|t|\le T$. On its left and horizontal
edges the exceptional reciprocal and the possible principal reciprocal
are each at most $1/\tau$ in absolute value. Thus the meromorphic
continuation
\[
        G_\chi(s)=\frac{\mathbf1_{\chi=1}}{s-1}
                                  -\frac{F'_\chi(s)}{F_\chi(s)}
\]
has absolute value at most $\mathcal M$ there. Its only poles in
the rectangle are the principal pole at one, of residue $1$ when
present, and the listed zeros $\beta>\sigma$, of residue $-1$.
The simplicity in \texttt{ExcZ} is essential for this last coefficient.

Apply the residue theorem to $G_\chi(s)u^s/s$. The two horizontal
edges have $|s|\ge T$, bounded length and $u^{\operatorname{Re}s}
\le u^a=eu\le3u$; their combined contribution is at most
$18\mathcal M u/T$. On the left edge $|s|\ge\sigma\ge1/2$,
giving at most $4T\mathcal M u^{1-2\tau}$.
If an exceptional zero lies to the left of this edge, inserting its
term in the final formula costs at most $2u^{1-2\tau}$, since
$\beta\ge1/2$ and $\beta\le\sigma$.
Consequently the Perron integral $J$ in (QR.41), with
$b=\lambda_\chi$, satisfies exactly the convenient source bound
\[
 \left|J-\mathbf1_{\chi=1}u+\sum_{\beta\in Z}\frac{u^\beta}{\beta}\right|
 \le18\mathcal M u/T+4T\mathcal M u^{1-2\tau}
                                      +2u^{1-2\tau}.       \tag{QR.43}
\]
The residue $1$ of the principal logarithmic derivative explains
the main term $u$; the residue $\rho$ of $L_1$ itself does not
multiply that main term.

\paragraph{QR22: choose a height uniformly in the conductor.}
Suppose $\log\mathcal B\le\sqrt{\log u}$ and put
$W=\sqrt{\log u}$, $\eta=c_0/32$. For $\eta W>1$, choose
$T=e^{\eta W}$. Then $2\le T\le u$, $W\ge1$, and
$\ell(T)\le3W$. Hence
\[
 \mathcal M\le K W^4,\quad K=9C(c_0)+48/c_0,
       \qquad u^{1-2\tau}\le u/T^2.                      \tag{QR.44}
\]
The last inequality follows from
$2\tau\log u=c_0W^2/(4\ell(T))\ge c_0W/12
\ge2\eta W$.
Adding (QR.41) and (QR.43) therefore bounds the total error by
$(2P+22K+2)W^4u/T$. The elementary exponential-series inequality
\[
 W^4e^{-\eta W}\le24(2/\eta)^4e^{-\eta W/2}
\]
absorbs this polynomial factor.

For the complementary range $\eta W\le1$, the elementary
Chebyshev estimate proved by the binomial argument in QW10 gives
$\sum_{n\le u}\Lambda(n)\le C_\psi u$ with an absolute
$C_\psi$. Thus $|\psi_\chi(u)|\le2C_\psi u$; the principal
term is at most $u$ and the exceptional term at most $2u$.
Since $e^{-\eta W/2}\ge1/2$, the same error estimate follows
after increasing its constant. One permissible uniform choice is
\[
 c_7=c_0/64,\qquad
 C_7=2(2C_\psi+3)+(2P+22K+2)24(2/\eta)^4.
\]
We have proved, for every family, character and admissible $Z$,
\[
 \left|\psi_\chi(u)-\mathbf1_{\chi=1}u
                         +\sum_{\beta\in Z}\frac{u^\beta}{\beta}\right|
       \le C_7u e^{-c_7\sqrt{\log u}},
       \quad u\ge2,\quad\log\mathcal B\le\sqrt{\log u}.  \tag{QR.45}
\]
This supplies the literal existential uniform explicit formula. The
human proof uses a coarse absolute Chebyshev constant in the small
range; the source's sharper value is not needed for its conclusion.

\paragraph{QR23: select a single family-wide exception and bound its interval loss.}
Choose once
\[
 0<c_0\le\min(1/20,c_{\rm zf},c_{\rm un}).
\]
For each character let $Z_\chi$ consist of its real zeros in
$[1-c_0/\ell(0),1)$. QR17 and QR18 say these sets satisfy
\texttt{ExcZ}, contain simple zeros, and have total cardinality
at most one across the entire character group. Each such character
is real. QR14 therefore supplies the common separation
\[
       1-\beta\ge\Delta_{\mathcal B}:=c_9\mathcal B^{-46}
           \quad(\beta\in Z_\chi),\qquad0<\Delta_{\mathcal B}\le1.
                                                               \tag{QR.46}
\]
Every such $\beta$ is at least $3/4$.
For $X\ge e$, $0\le h\le1$ and $Y=X(1+h)$, the exact difference
of its two residue terms is
\[
 \frac{Y^\beta-X^\beta}{\beta}
       =\int_X^Y t^{\beta-1}\,dt
       \le hX X^{\beta-1}
       \le hX e^{-\Delta_{\mathcal B}}
       \le hX(1-\Delta_{\mathcal B}/2).                   \tag{QR.47}
\]
The last inequality follows, for $0\le v\le1$, from convexity:
$e^{-v}\le1-(1-e^{-1})v\le1-v/2$.
If no exception exists the loss is zero, also bounded by the
nonnegative right side of (QR.47). Thus one common main-term
margin remains after summing over all characters; summing a separate
worst-case exceptional loss for every character would have destroyed
this argument.

\paragraph{QR24: obtain the full weighted short-interval lower bound.}
Take $A_0=46$ and the constants $c_9,c_7,C_7$ just proved. Choose
\[
 c=\min\left(1,\frac{c_7}{2(A_0+2)},
                          \frac{c_7c_9}{24C_7}\right)>0.  \tag{QR.48}
\]
Suppose $X\ge1$, $\log\mathcal B\le c\sqrt{\log X}$ and
$e^{-c\sqrt{\log X}}\le h\le1$. Write $W=\sqrt{\log X}$.
Since $\mathcal B\ge4$, one has $cW\ge\log\mathcal B\ge1$,
so $X\ge e$, and (QR.45) applies at both $X$ and $Y=X(1+h)$.
Character orthogonality gives the exact identity
\[
       \sum_\chi\psi_\chi(u)=|G|\sum_{n\le u}\lambda_0(n).
                                                               \tag{QR.49}
\]
The sum of the two errors in (QR.45), over the $|G|\le\mathcal B$
characters, is at most $3\mathcal B C_7X e^{-c_7W}$.
It is bounded by $hX\Delta_{\mathcal B}/4$. Indeed the ratio
of these two positive quantities is at most
\[
 \frac{12C_7}{c_9}
       e^{((A_0+2)c-c_7)W}
 \le\frac{12C_7}{c_9}e^{-c_7W/2}\le1.                   \tag{QR.50}
\]
The first step uses $\mathcal B\le e^{cW}$ and $h\ge e^{-cW}$;
the second uses (QR.48). For the last, $W\ge1/c$ and
$c\le c_7c_9/(24C_7)$ imply
$c_7W/2\ge12C_7/c_9$, and $e^t\ge t$ completes the bound.

Subtract (QR.49) at $X$ from its value at $Y$, use the single
exception estimate (QR.47), and then (QR.50). For the actual floor
interval sum $S=\sum_{\lfloor X\rfloor<n\le\lfloor Y\rfloor}
\lambda_0(n)\ge0$ this proves
\[
 |G|S\ge hX\Delta_{\mathcal B}/4,\qquad
 S\ge\frac{c_9hX}{4\mathcal B^{47}}.                     \tag{QR.51}
\]
This is the complete \texttt{trivialClass\_shortInterval} supplier
for every \texttt{LFam}, with the permissible concrete exponent
$47$. Its constants are chosen before the family, $X$ and $h$.
The lower bound retains the source restrictions on the conductor
and the relative interval length; it makes no assertion for arbitrary
short intervals.

\paragraph{QR25: close the arithmetic prime suppliers used by the reconstruction.}
Insert QR9's actual family, with $\mathcal B=4dM^2$ and norm-one
prime support, into (QR.51). Both inputs formerly retained in QW11
are now supplied. The complete prime-power removal in QW10--QW11
then gives a fixed $c_{\rm SN}>0$, chosen no larger than $c$, such
that
\[
 \#\{p\in(X,X(1+h)]:p\text{ prime},\ \mathrm{NormOne}(d,M,p)\}
       \ge\frac{c_9hX}{32(4dM^2)^{47}\log X},             \tag{QR.52}
\]
whenever $d$ is prime, $d\equiv3\pmod4$, $d\ge100$, $M\ge3$,
$X\ge1$, $\log(4dM^2)\le c_{\rm SN}\sqrt{\log X}$ and
$e^{-c_{\rm SN}\sqrt{\log X}}\le h\le1$.
For example use QW11 with its exponent $A=47$ and coefficient
$e=c_9/4$ to choose $c_{\rm SN}$ explicitly in terms of the
absolute Chebyshev remainder constant. This proves the exact
\texttt{ShortNorm} assertion, including the floor endpoints.

QW12's uniform scale comparison now proves the literal
\texttt{levelPrimeN} and its actual-host specialization: after the
single threshold depending only on the fixed source constants, one
finds a prime $T$ with
$T_s<T\le(1+1/V)T_s$ and
$\mathrm{NormOne}(d,N(g_0g_1)^2,T)$ for every admissible later
endpoint and datum. QW13 similarly gives the original
$Y^{9/10}$ lower bound in the dyadic norm-prime interval.
These conclusions no longer retain \texttt{ExistsLFamRayStatement}
or \texttt{TrivialClassShortIntervalStatement} as unproved inputs.
Their original \texttt{ModAt}, host and rank inequalities remain
the explicit hypotheses of the geometric applications. This supply
of norm primes does not itself choose compatible nonempty local
stages, simultaneous \texttt{PairEx} data, or the centered prime-sum
objects required later in the RH and Dirichlet-GRH branches.
